
\section{Seventh Sunday after Trinity: The True Banquet}


\begin{quote}

Luke 14:12–15. And he said also to the one who had invited him: When you make a dinner or a supper, do not invite your friends, your brothers, your relatives, or rich neighbors, lest they also invite you again and you receive repayment. But when you make a feast, invite the poor, the crippled, the lame, the blind. Then you shall be blessed; for they have nothing with which to repay you, but it shall be repaid to you at the resurrection of the righteous. But when one of those who sat with him at table heard this, he said to him: Blessed is the one who eats bread in the kingdom of God.

\end{quote}

\bigskip


According to this account and the one that precedes it, Jesus is a remarkable guest. He was in the house of a Pharisee and was sitting at the table with him. There he took the occasion to speak grave words both to the guests and to the host. To the guests he spoke about taking the lowest place and humbling themselves. To the host he spoke about whom he ought to invite when he gives a feast.

There is a sharp edge and sober seriousness in both of these table speeches. The guests were likely among those who sought honor from people rather than the honor that comes from the one God. And the host, who was a Pharisee, probably cared more for the reward he could receive at once than for the reward that lies hidden with the heavenly Father and will be revealed only on the day of Jesus Christ in the resurrection of the righteous.

The host invited friends and brothers and relatives and wealthy neighbors who could give a feast in return; but he did not think of making for himself such friends by means of unrighteous mammon as might receive him into the eternal dwellings.

Therefore the intention of Jesus Christ is not merely to give rules about what kinds of social gatherings may or may not be suitable within the congregation of God. His intention is to rebuke that whole form of religion and morality, that entire outlook on life and manner of living, which strives after earthly, temporal, and perishable advantages and distinctions while forgetting the eternal responsibility and the eternal and imperishable glory.

Large gatherings or small gatherings that pass, as it were, from one household to another in turn, and in which a brief earthly enjoyment is all that is sought and given, are not what Jesus has in mind. He always has eternity in view. And when he sees how people squander precious hours of grace through fleeting pleasures and empty diversions, and how they move thoughtlessly toward death and eternity, he has compassion on such misery. For he sees that they move along the edge of the abyss, only a hair’s breadth from eternal judgment.

He sees a spirit that cannot endure. Therefore he warns against gatherings that have no worth for eternity, and he commends a wholly different kind of gathering—one ordered entirely in view of the demands of eternal responsibility.

To feed the rich while the hungry and poor are sent away empty‑handed—that is the spirit of the world and the wisdom of the world. It receives its reward in this world and its comfort in the present time. But it is a grave folly for the person whose eye is directed toward what is eternal. For in the judgment it counts for nothing that a person has enjoyed life with his friends if he has neglected to clothe the naked, feed the hungry, and visit the sick and the sorrowful.

Give feasts for the poor. Help those who cannot help themselves. This is the counsel of Jesus to those who wish to give the right kind of feast. To take earthly wealth and turn it into eternal treasure, into a heavenly treasure that moth and rust cannot destroy and that no thief’s hand can carry away—this is the wisdom of the heavenly household and the true spiritual economy.

And this happens in a very simple and trusting way: the poor person’s hand becomes your bank, and you dare to entrust your means there. Whoever does this for God’s sake has stored both principal and interest with the Lord, and it will be returned to him in the resurrection of the righteous.

How can this be, and how does it hold together? If it could be explained to human reason, and if it were taken in a merely earthly way, it would become nothing more than a new kind of Phariseeism that imagines salvation can be purchased by almsgiving. That is not what Jesus intends.

It is not a bargain. It is not clever calculation. It is something entirely different: the boldness of faith and the self‑giving of love. For even if someone were to give all his possessions to the poor but lacked love, he could not—with all his wealth, even if it amounted to millions upon millions—buy even a moment of blessedness with God.

Many people, as death approaches, have given away everything they own to the poor—not out of love for their fellow human beings but out of fear of the pains of hell. Does God sell salvation for such a price? 

Certainly not.

There is only one ransom that delivers from death and hell and obtains salvation: the blood of Jesus, poured out for poor sinners. Yet the one in whom the love of Christ has become new life, so that he loves his neighbor and therefore helps him in his need—only such a person possesses the disposition that receives salvation from God in the love of God.

Faith in God and love toward brothers are necessary if one is to give the true feast for the poor, the crippled, the lame, and the blind. What is required is the faith that constantly recognizes the vanity of earthly things and trusts in the eternal, imperishable, and heavenly things.

It is the faith that turns its back on the world and fixes its gaze on the heavenly kingdom and on the misery here below. When the heart is loosened from earthly things because it grows together with the heavenly, then room is made for the right kind of feast. Then one learns the divine art of stewardship.

For it consists in sharing the perishable goods of the world with brothers who suffer want, because human hearts are gladdened by it. To give a suffering person a moment of joy, to bring true gladness to the heart, to relieve distress, to remove bitterness and grief from a burdened soul, to wipe the tears of someone who weeps in the love of Jesus Christ—this is what pleases God.

In this way a person gathers an imperishable treasure in heaven. Those who have been helped in this way will receive their helper and friend joyful and radiant in the eternal dwellings—provided they themselves have attained salvation with God.

And if someone whom you have helped in self‑giving love does not share in salvation, do not fear that the reward is lost. The Lord has seen in secret and will repay openly.

Thus the remarkable table discourse of Jesus in the Pharisee’s house becomes a blessed guide for those who sincerely wish to use the things of the earth rightly.

Blessed is the one who dares to order his household according to the instruction of Jesus. Such a person will lack nothing, neither on earth nor in heaven.

But woe to the one who guards earthly possessions so carefully that he does not dare to place them at interest with God.

We think so long about old age and about providing for our family that we often forget to think about the eternity that stands before us. We are so careful to secure a comfortable livelihood on earth that we forget the eternal punishment, the hunger that is never satisfied, and the thirst that is never quenched.

May the Lord teach us the true heavenly wisdom in the use of our earthly goods, so that we too may one day sit at the table in the kingdom of God.

For blessed is the one who may enjoy the glory of God forever and be satisfied with the sight of his face.